\chapter{Verification: From Free Electrons to Relativistic Gold}
\label{ch:verify}

Build and test in stages; each stage has an exact or well-established
target. Do not proceed to a stage until the previous one passes.

\paragraph{About the reference numbers in this chapter.} Every formula,
table and algorithm in this document has been exercised by an independent
reference implementation that follows the text literally (plain NumPy,
double precision, no symmetry beyond time reversal). Numbers quoted
``from the reference calculation'' come from it. They are not
experimental or literature values; they are what \emph{these equations}
give with \emph{these parameters}, and a correct implementation must
reproduce them to the stated digits. Where a number depends on
discretization choices (FFT grid size, mesh), those are stated.

\section{Stage 0: components in isolation}
\begin{enumerate}[leftmargin=2em]
\item \textbf{Basis}: plane-wave counts per $\kk$ close to
$\Omega(2E_{\mathrm{cut}})^{3/2}/6\pi^2$ (gold cell:
$a = 7.70668$\,Bohr, $\Omega = 114.4306$\,Bohr$^3$, so about $1380$ at
$E_{\mathrm{cut}} = 40$\,Ha).
\item \textbf{FFT round trip}: sphere $\to$ box $\to$ IFFT $\to$ FFT $\to$
gather returns the input to machine precision; Parseval's theorem holds.
\item \textbf{XC values}: the anchors of Chapter~\ref{ch:xc}
($\varepsilon_x, \varepsilon_c$ at $r_s = 2, 3$ to $10^{-5}$; the PBE
point to $10^{-6}$); AD-generated $\vxc$ matches the analytic derivative.
\item \textbf{Ewald}: independence of $\eta$ to $10^{-10}$. Quantitative
target for the gold cell ($Z_{\mathrm{ion}} = 11$,
$a = 4.0782$\,\AA, one atom): $E_{\mathrm{Ewald}} = -35.992686$\,Ha,
equivalent to the FCC Madelung constant,
$E_{\mathrm{Ewald}} = -0.8958736\, Z_{\mathrm{ion}}^2 / r_{\mathrm{ws}}$
with $r_{\mathrm{ws}} = (3\Omega/4\pi)^{1/3}$.
\item \textbf{Hermiticity}: $\braket{\phi|\hat H \psi} =
\overline{\braket{\psi | \hat H \phi}}$ for random spinors, with SOC on ---
catches conjugation bugs in projectors and $\mathbf L\cdot\mathbf S$.
\item \textbf{$\mathbf L\cdot\mathbf S$ blocks}: the three checks of
Section~\ref{sec:lsmatrix}, in the real-harmonic basis you actually use;
then the $10\times10$ crystal-field model of Chapter~\ref{ch:soc} must give
$\{-1, +1, +2.5\}\,\xi$ with degeneracies $4, 2, 4$ for $\Delta = 2\xi$.
\item \textbf{Symmetry}: $[\hat P,\hat H]=0$, $[\Theta,\hat H]=0$,
$\Theta^2=-1$ on random spinors at the eight TRIM
(Section~\ref{sec:symmetry}).
\item \textbf{Pseudopotential pieces}: the closed forms
\eqref{eq:hghlocG}--\eqref{eq:hghlocpoly} and \eqref{eq:hghprojG} against
brute-force radial quadrature of Eqs.~\eqref{eq:hghloc} and
\eqref{eq:besseltransform}; projector normalization
$\int p_{i\ell}^2 r^2 \dd r = 1$; and $\alpha$ of Eq.~\eqref{eq:alphaZ}:
$\alpha = 61.595082$ for the LDA parameters of Table~\ref{tab:auhgh}.
\item \textbf{Smearing}: $\sum_{n\kk} w_\kk f_{n\kk} = N_e$ after the
Fermi search; $\dd s/\dd x = x\,\dd f/\dd x$ numerically for each scheme
(this is the identity that makes $F$ variational).
\item \textbf{Eigensolver}: for a small cutoff, build the dense matrix by
applying \textsc{ApplyH} to unit vectors and compare \texttt{eigh} with
Algorithm~\ref{alg:lobpcg} --- after 30 iterations \emph{and} after 300.
\end{enumerate}

\section{Stage 1: exactly solvable systems}
\begin{enumerate}[leftmargin=2em]
\item \textbf{Empty lattice} ($V = 0$): eigenvalues must equal
$\tfrac12|\kk+\GG|^2$ exactly; band structure is the folded free-electron
parabola. Tests kinetic term, eigensolver, $\kk$ bookkeeping.
\item \textbf{Single cosine potential}: $V(\rr) = 2V_0\cos(\GG_1\cdot\rr)$
with $\GG_1 = \mathbf b_1$. At $\kk = \tfrac12\mathbf b_1$ (an $L$ point)
the two lowest levels must be $\tfrac18|\mathbf b_1|^2 \pm V_0$ up to
corrections of order $V_0^2/|\mathbf b_1|^2$, i.e.\ a gap
$\approx 2|V_0|$, in agreement with Section~\ref{sec:nfe} and with a dense
diagonalization; and their parities (Eq.~\eqref{eq:parityperm}) must be
even-below-odd for $V_0<0$ and the reverse for $V_0>0$.
\item \textbf{A screened ion in a big box.} This is the decisive test of
the pseudopotential code, SOC included, against an oracle that shares no
code with it. Put one gold ion in a large FCC cell ($a = 24$\,Bohr),
switch XC off, and use as the only other potential the Hartree potential
of a \emph{fixed} Gaussian charge of width $w = 2$\,Bohr,
$n(\GG) = (Z_{\mathrm{ion}}/\Omega)\,\ee^{-G^2w^2/2}$ --- the SCF starting
density with a different width. The total potential is then short-ranged
and spherically symmetric about the atom,
\[
V(r) = V_{\mathrm{loc}}(r) + \frac{Z_{\mathrm{ion}}}{r}\,
\mathrm{erf}\!\Big(\frac{r}{\sqrt2\,w}\Big) \;+\; \hat V_{\mathrm{NL}},
\]
so its bound states are those of a one-dimensional radial problem for each
$(\ell, j)$. Solve that problem independently: on a uniform radial grid
$r_i = ih$, the function $u(r) = r\,R(r)$ obeys a matrix eigenproblem with
\begin{multline*}
H_{ii'} = -\frac{\delta_{i,i'+1} - 2\delta_{ii'} + \delta_{i,i'-1}}{2h^2}
+ \Big[\frac{\ell(\ell+1)}{2r_i^2} + V(r_i)\Big]\delta_{ii'}\\
+ h\sum_{a,b=1}^{2} r_i\,p_{a\ell}(r_i)\;\big(h^\ell_{ab} + \lambda_{\ell j}\,k^\ell_{ab}\big)\;
r_{i'}\,p_{b\ell}(r_{i'}),
\end{multline*}
where $\lambda_{\ell j} = \ell/2$ for $j = \ell+\tfrac12$ and
$-(\ell+1)/2$ for $j = \ell-\tfrac12$ (Eq.~\eqref{eq:ls_eigen}) --- thirty
lines of code, one dense \texttt{eigh} per channel ($h = 0.01$,
$r_{\max} = 30$). Reference values (extrapolated to $h\to0$):
\begin{center}
\begin{tabular}{@{}lccccc@{}}
\toprule
 & $s_{1/2}$ & $p_{1/2}$ & $p_{3/2}$ & $d_{3/2}$ & $d_{5/2}$\\
\midrule
$\varepsilon$ (Ha) & $-1.01843$ & $-0.60281$ & $-0.59335$ & $-1.99901$ & $-1.89459$\\
degeneracy in the box & 2 & 2 & 4 & 4 & 6\\
\bottomrule
\end{tabular}
\end{center}
The plane-wave eigenvalues at $\Gamma$ must equal these \emph{plus}
$2\pi Z_{\mathrm{ion}}w^2/\Omega$, with the degeneracies shown. (The shift
is the price of setting $\vH(\GG{=}0)=0$: the plane-wave potential differs
from the true one by exactly this constant --- a useful exercise in
Section~\ref{sec:g0}.) The reference plane-wave calculation ($E_{\mathrm{cut}} = 40$\,Ha,
$41\,655$ plane waves, FFT grid $98^3$, shift $0.079994$\,Ha) gives
$-1.01843$, $-0.60281$, $-0.59335$, $-1.99900$, $-1.89458$\,Ha: agreement
to $10^{-5}$\,Ha in every channel. The $d_{5/2}$--$d_{3/2}$
splitting, $0.1044$\,Ha with $j = \tfrac52$ \emph{above} $j=\tfrac32$,
fixes the overall sign of the $k^\ell_{ij}$ term.
\item \textbf{The self-consistent atom (oracle only).} Adding a Hartree and
LDA potential to the radial solver and iterating gives the neutral
pseudo-atom. For the configuration $5d_{3/2}^{4}\,5d_{5/2}^{6}\,6s^1$
(spherical, not spin-polarized, LDA parameters):
$\varepsilon_{6s} = -0.2241$, $\varepsilon_{5d_{3/2}} = -0.2971$,
$\varepsilon_{5d_{5/2}} = -0.2405$\,Ha, hence a spin--orbit splitting of
the $5d$ level of $1.54$\,eV --- the number quoted in
Chapter~\ref{ch:soc}, here derived from the twenty numbers of
Table~\ref{tab:auhgh}. (The unoccupied $6p$ levels come out at $-0.0486$
and $-0.0278$\,Ha, split by $0.57$\,eV with $j=\tfrac12$ lower.)
\end{enumerate}

\section{Stage 2: self-consistency on easy targets}
\begin{enumerate}[leftmargin=2em]
\item \textbf{Jellium in a cell}: no ions, uniform positive background
(i.e.\ just $\vH(\GG{=}0) = 0$ and no ionic potential), $N_e$ electrons.
The SCF must return the uniform density $N_e/\Omega$, eigenvalues
$\tfrac12|\kk+\GG|^2 + \vxc(n)$, and the analytic HEG energy per electron
$\tfrac{3}{10}k_F^2 + \varepsilon_{\mathrm{xc}}(n)$ in the limit of a dense
mesh and small smearing.
\item \textbf{Silicon (bonus)}: if you want a soft interim target, HGH
silicon converges with $E_{\mathrm{cut}} \approx 12$\,Ha; indirect
gap $\approx 0.5$--$0.6$\,eV (the famous LDA/PBE underestimate), lattice
constant PBE $\approx 5.47$\,\AA. (Two atoms per cell: your structure
factors and projector phases get their first real test.)
\end{enumerate}

\section{Stage 3: gold, self-consistent --- regression numbers}
\label{sec:stage3}

Start small, and start with spin--orbit coupling switched off
($k^\ell = 0$): the three cases of Table~\ref{tab:smallcases} run in a
minute on a laptop CPU and each isolates one ingredient. Work through them
in order --- scalar-relativistic LDA first, then SOC, then PBE.

\begin{table}[H]
\centering\small
\begin{tabular}{@{}lccl@{}}
\toprule
Case & $F$ (Ha) & $\mu$ (Ha) & levels at $\Gamma$ in Ha (number of spinor states)\\
\midrule
LDA, SOC off & $-33.291161$ & $0.613238$ & $0.213490\,(2)$, $0.406739\,(6)$, $0.463560\,(4)$\\
LDA, SOC on & $-33.311413$ & $0.609888$ & $0.213402\,(2)$, $0.381623\,(4)$, $0.425863\,(2)$, $0.470840\,(4)$\\
PBE, SOC on & $-33.254853$ & $0.624314$ & $0.229208\,(2)$, $0.394565\,(4)$, $0.438491\,(2)$, $0.482719\,(4)$\\
\bottomrule
\end{tabular}
\caption{Small regression cases. Common settings: $q = 11$ parameters of
Table~\ref{tab:auhgh} matching the functional, $a = 4.0782$\,\AA{}
$= 7.706681$\,Bohr, $E_{\mathrm{cut}} = 20$\,Ha, FFT grid $24^3$ ($531$
plane waves at $\Gamma$), $\Gamma$-centered $4^3$ mesh ($36$ points after
time-reversal reduction; density accumulated as $n$ only, with
$\mathbf m = 0$, as Section~\ref{sec:symmetry} requires on such a mesh), Methfessel--Paxton order 1 with
$\sigma = 0.01$\,Ha, 20 spinor bands, SCF as in the driver of
Chapter~\ref{ch:algorithms}, converged to $\|R\| < 10^{-8}$ (10--12
iterations). The levels at $\Gamma$ are the eigenvalues of the converged
Hamiltonian; with SOC off the spinor code returns every level of a scalar
code twice (the $t_{2g}$ triplet as six states, the $e_g$ doublet as
four). The cutoff is far too low for converged physics --- these are
regression numbers.}
\label{tab:smallcases}
\end{table}

A correct implementation reproduces $F$, $\mu$ and the levels to the last
digit shown: both the NumPy reference and an independent JAX program
written from this text do. If you are off, the size of the discrepancy
tells you where to look:
\begin{itemize}[leftmargin=2em]
\item \emph{$10^{-6}$--$10^{-5}$\,Ha in $F$}: rounded constants (PBE
$\beta,\mu$; pseudopotential digits; PW92 $A$) or a different FFT grid.
\item \emph{$10^{-3}$\,Ha or more in $F$ with correct eigenvalue
differences}: a constant --- $\alpha$ of Eq.~\eqref{eq:alphaZ}, the Ewald
energy, the entropy term, or double counting in Eq.~\eqref{eq:ftot}.
\item \emph{Wrong degeneracies at $\Gamma$} (anything other than $2,6,4$
without and $2,4,2,4$ with SOC): projectors, real harmonics or the
$\mathbf L\cdot\mathbf S$ blocks; go back to Stage~0 and Stage~1.3.
\item \emph{SOC-off case right, SOC-on wrong}: the $k^\ell_{ij}$ values or
the assembly of $D$.
\item \emph{LDA right, PBE wrong}: the gradient or divergence step of
Eq.~\eqref{eq:vxcgga2}; test the PBE point of Chapter~\ref{ch:xc} first.
\end{itemize}

Two larger cases with spin--orbit coupling on, for the production
settings:
\begin{center}\small
\begin{tabular}{@{}lll@{}}
\toprule
Setting & Case A & Case B\\
\midrule
functional, parameters & LDA, $q=11$ & LDA, $q=11$\\
$a$ & $7.706681$\,Bohr & $7.706681$\,Bohr\\
$E_{\mathrm{cut}}$; FFT grid & $25$\,Ha; $27^3$ & $40$\,Ha; $35^3$\\
mesh (time-reversal reduced) & $4^3$, Eq.~\eqref{eq:mpmesh}: 32 points & $8^3$, $\Gamma$-centered: 260 points\\
smearing & MP1, $\sigma=0.01$\,Ha & MP1, $\sigma=0.01$\,Ha\\
spinor bands & 20 & 20\\
\midrule
$F$ (Ha) & $-33.324309$ & $-33.325815$\\
$\mu$ (Ha) & $0.58409$ & $0.58322$\\
SCF iterations to $\|R\|<10^{-7}$ & 10 & 10\\
\bottomrule
\end{tabular}
\end{center}
(Throughout: $a_0 = 0.529177211$\,\AA; PW92 constants of
Table~\ref{tab:pw92}. Case~A deliberately uses the shifted mesh of
Eq.~\eqref{eq:mpmesh} to exercise that code path; it is a regression
number, not a recommendation --- see Section~\ref{sec:bzintegration}.)
Case~B is the calculation behind
Table~\ref{tab:goldlevels}, Fig.~\ref{fig:bands} and
Chapter~\ref{ch:topology}.

\paragraph{Then the physics.} Lattice constant and bulk modulus from
equation-of-state fits (end of Chapter~\ref{ch:hamiltonian}; seven
lattice constants spanning $\pm4\%$, scalar-relativistic,
$E_{\mathrm{cut}} = 45$\,Ha, $\Gamma$-centered $8^3$ mesh, MP1
$\sigma = 0.01$\,Ha), and the position of the $d$ bands:
\begin{center}\small\setlength{\tabcolsep}{4pt}
\begin{tabular}{@{}lccccc@{}}
\toprule
 & \multicolumn{2}{c}{LDA} & \multicolumn{2}{c}{PBE} & Experiment\\
 & this potential & literature & this potential & literature & (room temp.)\\
\midrule
Lattice constant $a_0$ (\AA) & $4.093$ & $4.053$ & $4.200$ & $4.173$ & $4.078$\\
Bulk modulus $B_0$ (GPa) & $184$ & $195$ & $133$ & $143$ & $173$\\
$B_0'$ & $5.7$ & & $6.1$ & & \\
$d$-band top at $L$ (eV below $E_F$) & $1.63$ & & & $1.90$ & $2.3$\\
\bottomrule
\end{tabular}
\end{center}
``This potential'' is the reference calculation with the 11-electron
parameters of Table~\ref{tab:auhgh}; ``literature'' are
all-electron-quality values (Ambrosetti and Silvestrelli 2016 for $a_0$
and $B_0$; Rangel \emph{et al.}\ 2012 for the $d$ band, which is quoted
with spin--orbit coupling). Three lessons.
(i) The familiar statement ``LDA overbinds, PBE underbinds'' is visible in
both pairs of columns: PBE gives a lattice constant $0.11$--$0.12$\,\AA{}
larger and a bulk modulus about $50$\,GPa smaller than LDA. (ii) The
11-electron pseudopotential adds $0.04$\,\AA{} (LDA) and $0.03$\,\AA{} (PBE)
to $a_0$ and takes
about $10$\,GPa off $B_0$ compared with all-electron-quality
calculations --- enough to move the LDA lattice constant from slightly
below to slightly above experiment. This is a property of the
pseudopotential (no semicore states, no nonlinear core correction), not
of your code: reproduce the first and third columns. (iii) Spin--orbit
coupling changes these numbers little (Stage~4): at
$E_{\mathrm{cut}} = 35$\,Ha with a $6^3$ mesh, LDA gives
$4.095$\,\AA, $184$\,GPa without and $4.084$\,\AA, $190$\,GPa with SOC;
comparison with the first column also shows that the equation of state
is converged with respect to cutoff and mesh at the level of
$0.003$\,\AA{} and $1$\,GPa. The $d$-band position discrepancy versus
photoemission and optics is a known limitation of LDA and GGA ---
reproduce the DFT number, not experiment. The SCF for the metal must converge in
about ten iterations with Kerker--Pulay mixing; watch the difference
between the two energy expressions \eqref{eq:fdirect} and \eqref{eq:ftot}
close quadratically. A typical history (reference calculation,
$E_{\mathrm{cut}} = 40$\,Ha, $8^3$ mesh, SOC on):
\begin{center}\small
\begin{tabular}{@{}cccc@{}}
\toprule
iteration & $\|R\|$ & $F$ (Ha), Eq.~\eqref{eq:ftot} & Eq.~\eqref{eq:fdirect} $-$ Eq.~\eqref{eq:ftot}\\
\midrule
1 & $6.1\times10^{-1}$ & $-33.541764$ & $4.2\times10^{-1}$\\
2 & $2.3\times10^{-1}$ & $-33.330984$ & $1.3\times10^{-2}$\\
3 & $5.4\times10^{-2}$ & $-33.325191$ & $-5.9\times10^{-4}$\\
4 & $1.8\times10^{-3}$ & $-33.325816$ & $2.1\times10^{-6}$\\
5 & $5.9\times10^{-4}$ & $-33.325816$ & $6.0\times10^{-7}$\\
6 & $3.4\times10^{-5}$ & $-33.325815$ & $-6.2\times10^{-10}$\\
8 & $2.0\times10^{-6}$ & $-33.325815$ & $7\times10^{-12}$\\
10 & $4.7\times10^{-8}$ & $-33.325815$ & $-4\times10^{-13}$\\
\bottomrule
\end{tabular}
\end{center}

\section{Stage 4: gold with spin--orbit coupling --- the physics}
\begin{enumerate}[leftmargin=2em]
\item \textbf{Kramers pairs}: every eigenvalue doubly degenerate at
every $\kk$ (inversion $\times$ time reversal), to the accuracy of the
eigensolver (the reference calculation finds $4\times10^{-11}$\,Ha on the
SCF mesh after three LOBPCG iterations per step).
\item \textbf{Cubic degeneracies}: with a $\Gamma$-centered mesh, the
quartets at $\Gamma$ are exact and all symmetry-equivalent $\kk$-points
agree to solver accuracy.
\item \textbf{Level scheme at the high-symmetry points},
Table~\ref{tab:goldlevels}: without SOC two $d$ levels at $\Gamma$ (a
$t_{2g}$ triplet and an $e_g$ doublet, each $\times2$ for spin); with SOC
three, with degeneracies $4, 2, 4$ in ascending order. Published
fully-relativistic band structures of Au agree with these values to
within the differences between pseudopotentials and functionals
($\sim0.1$\,eV).
\item \textbf{Energetics barely move}: switching SOC on lowers the
free energy by about $0.02$\,Ha, reduces $a_0$ by $0.011$\,\AA{} and
raises $B_0$ by $6$\,GPa (reference equation-of-state fits, LDA,
$E_{\mathrm{cut}} = 35$\,Ha, $6^3$ mesh: $4.095\to4.084$\,\AA,
$184\to190$\,GPa) --- a check that SOC entered as a (strong) band
effect, not as an energy-scale bug.
\item \textbf{Switch test}: SOC on/off touches only the $k^{\ell}$ blocks;
with $k^{\ell} = 0$ the spinor code must return every eigenvalue exactly
twice (and, if you also wrote a scalar code with two electrons per
orbital, reproduce its eigenvalues and free energy to solver precision).
\end{enumerate}

\begin{table}[H]
\centering\small
\begin{tabular}{@{}lll@{}}
\toprule
Point & SOC & Levels (eV relative to $E_F$)\\
\midrule
$\Gamma$ & off & $-10.06$,\ \ $-4.82\,(\times3)$,\ \ $-3.28\,(\times2)$,\ \ $15.97$\\
         & on  & $-10.06$,\ \ $-5.45\,(\times2)$,\ \ $-4.24$,\ \ $-3.02\,(\times2)$,\ \ $15.97$\\
\midrule
$X$ & off & $-7.20$,\ \ $-6.88$,\ \ $-2.09$,\ \ $-1.70\,(\times2)$,\ \ $0.96$,\ \ $5.76$\\
    & on  & $-7.32$,\ \ $-7.03$,\ \ $-2.31$,\ \ $-2.09$,\ \ $-1.05$,\ \ $0.71$,\ \ $5.78$\\
\midrule
$L$ & off & $-7.47$,\ \ $-4.86\,(\times2)$,\ \ $-2.03\,(\times2)$,\ \ $-1.06$,\ \ $3.07$\\
    & on  & $-7.56$,\ \ $-5.57$,\ \ $-4.38$,\ \ $-2.36$,\ \ $-1.63$,\ \ $-1.22$,\ \ $3.13$\\
\midrule
$W$ & off & $-6.10$,\ \ $-5.28\,(\times2)$,\ \ $-3.42$,\ \ $-1.70$,\ \ $4.71$,\ \ $6.40\,(\times2)$\\
    & on  & $-6.31$,\ \ $-5.73$,\ \ $-5.15$,\ \ $-3.25$,\ \ $-1.63$,\ \ $4.70$,\ \ $5.60$,\ \ $6.91$\\
\midrule
$K$ & off & $-6.59$,\ \ $-5.97$,\ \ $-3.95$,\ \ $-2.98$,\ \ $-2.16$,\ \ $3.85$,\ \ $4.73$\\
    & on  & $-6.75$,\ \ $-6.17$,\ \ $-4.06$,\ \ $-3.01$,\ \ $-1.95$,\ \ $3.76$,\ \ $4.63$\\
\bottomrule
\end{tabular}
\caption{Kohn--Sham levels of gold in eV relative to $E_F$ from the
reference calculation (LDA, $q=11$, $a = 4.0782$\,\AA,
$E_{\mathrm{cut}} = 40$\,Ha, $\Gamma$-centered $8^3$ mesh, MP1
$\sigma = 0.01$\,Ha; $\mu = 0.58322$\,Ha, $F = -33.325815$\,Ha). ``SOC
off'' levels were evaluated in the same self-consistent potential with
$k^\ell = 0$. Multiplicities count spatial states (SOC off) or Kramers
pairs (SOC on). Compare with Fig.~\ref{fig:bands}.}
\label{tab:goldlevels}
\end{table}

\section{Stage 5: topology}
\begin{enumerate}[leftmargin=2em]
\item The sanity checks of Chapter~\ref{ch:topology} (empty-lattice
parities, the cosine-potential inversion, $\Pi^2 = \mathbf 1$,
$[\hat P,\hat H] = 0$, all-even at $\Gamma$, SOC-off parities,
Wilson-loop pinning, robustness).
\item Target: $\delta_\Gamma = +1$, $\delta_X = -1$, $\delta_L = -1$,
hence $(\nu_0;\nu_1\nu_2\nu_3) = (1;000)$ for the lowest six Kramers
pairs, with a strictly positive direct gap between spinor bands 12 and 13
(Section~\ref{sec:directgap}).
\item Wilson loops: the end-point pattern of Fig.~\ref{fig:wcc} (one
Kramers pair at $\theta=\pi$ only on the line through $\Gamma$), and,
with sufficient resolution, $z(k_3{=}0) = 1$, $z(k_3{=}\tfrac12) = 0$,
hence the same $\nu_0 = 1$, $\nu_3 = 0$.
\end{enumerate}

\section{Stage 6: transport and optics}
\begin{enumerate}[leftmargin=2em]
\item \textbf{Velocity operator}: the Hellmann--Feynman test of
Section~\ref{sec:velocity} at a generic $\kk$ (the reference uses
fractional $\kk = (0.13, 0.27, 0.41)$): finite-difference band slopes and
pair-averaged diagonal velocity elements agree to $10^{-7}$\,a.u.; the
velocity matrices are Hermitian; dropping the nonlocal term changes the
velocities by $0.1$--$0.3$\,a.u.
\item \textbf{Fermi-surface numbers} (fixed potential of Case~B,
irreducible wedge of a $\Gamma$-centered $40^3$ mesh, Gaussian width
$0.005$\,Ha): $g(E_F) = \GFb$ states/eV per atom,
$\hbar\omega_p = \WPb$\,eV, r.m.s.\ Fermi velocity
$\VFb\times10^{6}$\,m/s, effective number of carriers \NEFF. Expect
agreement to a few percent, not to all digits: these are Fermi-surface
integrals and depend on mesh and broadening.
\item \textbf{Colour chain}: $R\equiv1$ gives sRGB $(255,255,255)$;
$R\equiv0.5$ gives $(188,188,188)$; a mirror that reflects only for
$\lambda > 600$\,nm gives $(234,0,0)$; one that reflects only for
$\lambda < 480$\,nm gives $(67,0,255)$ (strict inequalities on the 1\,nm
grid; with $\ge$ and $\le$ the values become $(237,0,0)$ and
$(65,0,255)$).
\item \textbf{Spectrum}: interband $\varepsilon_2$ below 1 up to
$\approx1.2$\,eV and above 2 from $\ONSET$\,eV, reflectivity above $0.9$ in
the near infrared and falling through $0.7$ at $\EDGE$\,eV
(Fig.~\ref{fig:optics}); Kramers--Kronig consistency of
Eq.~\eqref{eq:epsinter}; sum rule of Chapter~\ref{ch:colour}.
\end{enumerate}

\section{Stage 7: convergence studies (the habits of a professional)}
For each observable, produce convergence curves in
(i) $E_{\mathrm{cut}}$ at fixed mesh, (ii) $\kk$-mesh at fixed
$E_{\mathrm{cut}}$, (iii) smearing $\sigma \to 0$ jointly with the mesh.
Quote every number in your write-up with the settings that converge it to a
stated tolerance. This discipline --- not any single number --- is what makes
a DFT calculation trustworthy.

\section*{Further reading}
\begin{itemize}[leftmargin=2em]
\item R.\ M.\ Martin, \emph{Electronic Structure: Basic Theory and Practical
Methods} (CUP) --- the encyclopedic companion to everything here.
\item M.\ C.\ Payne et al., Rev.\ Mod.\ Phys.\ \textbf{64}, 1045 (1992) ---
the classic plane-wave/iterative-methods review.
\item C.\ Hartwigsen, S.\ Goedecker, J.\ Hutter, Phys.\ Rev.\ B \textbf{58},
3641 (1998) --- the relativistic HGH pseudopotential tables (your gold
parameters, including SOC).
\item J.\ P.\ Perdew, K.\ Burke, M.\ Ernzerhof, Phys.\ Rev.\ Lett.\
\textbf{77}, 3865 (1996); J.\ P.\ Perdew, Y.\ Wang, Phys.\ Rev.\ B
\textbf{45}, 13244 (1992) --- PBE and PW92.
\item A.\ V.\ Knyazev, SIAM J.\ Sci.\ Comput.\ \textbf{23}, 517 (2001) ---
LOBPCG; U.\ Hetmaniuk, R.\ Lehoucq, J.\ Comput.\ Phys.\ \textbf{218}, 324
(2006), and J.\ A.\ Duersch et al., SIAM J.\ Sci.\ Comput.\ \textbf{40},
C655 (2018) --- why and how to keep its basis orthonormal.
\item G.\ Kresse, J.\ Furthm\"uller, Phys.\ Rev.\ B \textbf{54}, 11169
(1996) --- iterative diagonalization, Pulay mixing and Kerker
preconditioning for metals.
\item P.\ Pyykk\"o, Chem.\ Rev.\ \textbf{88}, 563 (1988) --- relativity in
chemistry; why gold is gold.
\item M.\ Krack, Theor.\ Chem.\ Acc.\ \textbf{114}, 145 (2005) --- the PBE
parameter sets of the same pseudopotential family.
\item D.\ Vanderbilt, \emph{Berry Phases in Electronic Structure Theory}
(CUP, 2018) --- the book-length version of Chapter~\ref{ch:topology}.
\item L.\ Fu, C.\ L.\ Kane, Phys.\ Rev.\ B \textbf{76}, 045302 (2007) ---
the parity formula for the $\mathbb{Z}_2$ invariant.
\item A.\ A.\ Soluyanov, D.\ Vanderbilt, Phys.\ Rev.\ B \textbf{83}, 235401
(2011) --- the Wilson-loop (Wannier charge center) method.
\item B.\ Yan et al., Nat.\ Commun.\ \textbf{6}, 10167 (2015) ---
topological states on the gold surface.
\item S.\ LaShell, B.\ A.\ McDougall, E.\ Jensen, Phys.\ Rev.\ Lett.\
\textbf{77}, 3419 (1996) --- the spin-split surface state of Au(111).
\item N.\ W.\ Ashcroft, N.\ D.\ Mermin, \emph{Solid State Physics},
chapters 1, 12 and 13 --- the Drude model and semiclassical transport
behind Chapter~\ref{ch:conductivity}.
\item P.\ B.\ Allen, Phys.\ Rev.\ B \textbf{17}, 3725 (1978) --- the
electron--phonon theory of the resistivity of metals
(Eq.~\eqref{eq:rhoT}); S.\ Y.\ Savrasov, D.\ Y.\ Savrasov, Phys.\ Rev.\ B
\textbf{54}, 16487 (1996) --- first-principles $\lambda_{\mathrm{tr}}$
and resistivities of eight metals; T.\ Gunst, A.\ Blom, K.\ Stokbro,
Proc.\ SISPAD 2018, p.~79 --- phonon-limited resistivity of Au, Ag, Cu.
\item T.\ Rangel et al., Phys.\ Rev.\ B \textbf{86}, 125125 (2012) ---
the band structure of gold from many-body perturbation theory, with a
table of measured band energies (Table~\ref{tab:dbandexp}).
\item A.\ Ambrosetti, P.\ L.\ Silvestrelli, Phys.\ Rev.\ B \textbf{94},
045124 (2016) --- lattice constants and bulk moduli of Cu, Ag, Au with
LDA, PBE and beyond.
\item P.\ B.\ Johnson, R.\ W.\ Christy, Phys.\ Rev.\ B \textbf{6}, 4370
(1972) --- the measured optical constants of the noble metals (the
tables used here are in \texttt{verification/reference\_data/}, from
the public-domain refractiveindex.info database); R.\ L.\ Olmon et al.,
Phys.\ Rev.\ B \textbf{86}, 235147 (2012) --- a modern remeasurement
for gold.
\item C.\ Wyman, P.-P.\ Sloan, P.\ Shirley, J.\ Comput.\ Graph.\ Tech.\
\textbf{2}(2), 1 (2013) --- the analytic colour-matching functions of
Eq.~\eqref{eq:cmf}.
\item The DFTK.jl and GPAW documentation --- two readable open-source
plane-wave codes to compare designs against (resist reading their source
until yours works).
\end{itemize}

\vspace{1em}
\begin{keybox}{You are ready}
Every equation your program needs is now on paper: the spinor Kohn--Sham
problem \eqref{eq:spinorks}, its plane-wave matrix elements
(Chapter~\ref{ch:hamiltonian}), the XC input (Chapter~\ref{ch:xc}), the
relativistic ionic data (Chapters~\ref{ch:pseudo}--\ref{ch:soc}), and the
solvers (Chapter~\ref{ch:algorithms}) in GPU-native form
(Chapter~\ref{ch:gpu}). Build the program in the order of this chapter:
components, empty lattice, one screened gold ion, scalar-relativistic
gold, the full spin--orbit band structure of the metal that is yellow
because of special relativity, its topological invariant, and finally
its Fermi-surface and optical properties.
\end{keybox}
